\chapter{Extended Details}
\label{chap:extended-details}

\section{Additional Robust SSL Related Work}
\label{sec:robust-ssl-additional}

The main text centers on \DeACL. Other robust SSL methods frame the setting: RoCL uses adversarial contrastive examples \parencite{kim2020adversarial}; ACL enforces adversarial consistency \parencite{jiang2020robust}; AdvCL adds pseudo-labels \parencite{fan2021does}; CLAE builds adversarial contrastive examples \parencite{ho2020CLAE}; and DYNACL schedules augmentations around the augmentation--robustness conflict \parencite{luo2023rethinking}.

These methods show robustness-aware SSL can work, often by changing pre-training. \DeACL separates large-scale SSL pre-training from robust fine-tuning. This thesis does not dismiss representation fine-tuning: it strongly helps against \EmbedAttack. It shows that representation-only objectives do not guarantee downstream-loss robustness.

\section{Depth Loss Details}
\label{sec:depth-loss-details}

Depth estimation combines scale-invariant pixel error with spatial gradient matching. Let $g_i=\log \tilde{d}_i-\log d_i$, where $\tilde{d}_i$ is predicted depth and $d_i$ is ground truth. The scale-invariant pixel loss is
\begin{equation}
    L_{\mathrm{pixel}}
    =
    \alpha
    \sqrt{
        \frac{1}{T}\sum_i g_i^2
        -
        \frac{\rho}{T^2}\left(\sum_i g_i\right)^2
    },
\end{equation}
with $\alpha=1$ and $\rho=0.85$ in the supplied setup. Multi-scale gradient loss penalizes log-depth residual-gradient differences:
\begin{equation}
    L_{\mathrm{grad}}=
    \frac{1}{n}\sum_k\sum_i
    \left|\nabla_x R_i^k\right|+
    \left|\nabla_y R_i^k\right|.
\end{equation}
The depth training and attack objective is
\begin{equation}
    L_{\mathrm{depth}}=\frac{1}{2}L_{\mathrm{grad}}+L_{\mathrm{pixel}}.
\end{equation}

\section{Hyperparameters}
\label{sec:hyperparameters}

\begin{table}[htbp]
    \centering
    \small
    \begin{tabularx}{\textwidth}{lX}
    \toprule
    Component & Hyperparameters \\
    \midrule
    \EmbedAttack & $\ell_\infty$ budget $\varepsilon=8/255$, 20 PGD steps, step size $2/255$, random uniform initialization in the $\varepsilon$ ball. \\
    Classification \PGD & $\ell_\infty$ budget $\varepsilon=8/255$, 20 PGD steps, step size $2/255$, cross-entropy objective. \\
    \SegPGD & Same budget and PGD schedule, with weighted pixel-wise cross-entropy over correct and incorrect pixels. \\
    \DepthPGD & Same budget and PGD schedule, with the depth-estimation loss in \Cref{sec:depth-loss-details}. \\
    Classification head & Adam, learning rate 0.5, batch size 16, 5 epochs, random horizontal flips. \\
    Segmentation head & AdamW, learning rate $10^{-4}$, batch size 16, weight decay 0.001, 50 epochs, random cropping, horizontal and vertical flips, sliding-window inference. \\
    Depth head & AdamW, learning rate $10^{-4}$, batch size 128, weight decay 0.01, 20 epochs. \\
    \DeACL fine-tuning & SGD with momentum 0.9, learning rate 0.05 with cosine schedule, 10 warmup epochs, 100 total epochs, batch size 128, $\gamma=2$, adversarial budget $4/255$, random crops and flips. \\
    \bottomrule
    \end{tabularx}
    \caption{Hyperparameters from the supplied SSL robustness resources.}
    \label{tab:hyperparameters}
\end{table}

\section{Dataset Notes}
\label{sec:dataset-notes}

CIFAR10 and CIFAR100 are small-image classification datasets for fast linear-probe tests. CIFAR10 has ten coarse classes; CIFAR100 has one hundred finer classes. Both show that high clean accuracy does not prevent task-specific PGD failure.

STL10 contains natural images used in representation learning. It gives \DeACL the largest task-specific PGD value. This does not overturn the conclusion, but suggests dataset distribution affects transfer from representation to classification robustness.

ADE20K, Cityscapes, and PASCAL VOC 2012 cover diverse scenes, urban driving, and object-centric images. Using all three avoids a one-dataset segmentation conclusion.

NYU-Depth v2 is an indoor benchmark with continuous outputs, testing whether perturbations damage geometry, not only labels.

\section{Attack Pseudocode}
\label{sec:attack-pseudocode}

All PGD-style attacks share one procedure; only losses differ.

\begin{enumerate}[wide]
    \item Choose a clean image $x$, target $y$, pipeline $F$, budget $\varepsilon$, step size $\alpha$, and steps $K$.
    \item Initialize $\delta_0$ uniformly in $[-\varepsilon,\varepsilon]$ and clip $x+\delta_0$ to the valid pixel range.
    \item For each step $k=0,\ldots,K-1$, compute the selected loss on $x+\delta_k$.
    \item Update $\delta_{k+1}$ in the gradient-sign direction.
    \item Project $\delta_{k+1}$ to $[-\varepsilon,\varepsilon]$ and clip the perturbed image.
    \item Evaluate the final adversarial image with the task metric.
\end{enumerate}

\EmbedAttack uses representation distance; classification \PGD uses cross-entropy; \SegPGD uses weighted pixel-wise cross-entropy; and \DepthPGD uses depth loss. The shared skeleton enables comparison under one perturbation budget.

\section{Notation Summary}
\label{sec:notation}

\begin{table}[htbp]
    \centering
    \small
    \begin{tabularx}{\textwidth}{lX}
    \toprule
    Symbol & Meaning \\
    \midrule
    $x$ & Clean input image. \\
    $x_{\mathrm{adv}}$ & Adversarially perturbed image. \\
    $\delta$ & Perturbation added to the clean image. \\
    $\varepsilon$ & Maximum $\ell_\infty$ perturbation budget. \\
    $f_\theta$ & Self-supervised vision encoder. \\
    $f_R$ & Robust encoder produced by \DeACL fine-tuning. \\
    $h_\phi^t$ & Downstream adaptor for task $t$. \\
    $L_{\mathrm{CE}}$ & Cross-entropy classification loss. \\
    $L_{\mathrm{seg}}$ & Segmentation attack loss. \\
    $L_{\mathrm{depth}}$ & Depth-estimation loss used by \DepthPGD. \\
    \bottomrule
    \end{tabularx}
    \caption{Notation used throughout the thesis.}
    \label{tab:notation}
\end{table}

\section{Additional Validity Threats}
\label{sec:validity-additional}

Several validity threats remain. The robust encoder is only available for \DINO ViT-B/16, limiting architectural conclusions. The lightweight heads may differ from stronger heads in clean metrics and attack surfaces. The attacks are digital white-box attacks; physical perturbations, distribution shift, corruptions, and black-box attacks are excluded.

Thesis expansion is not experimental expansion: the rewrite explains supplied experiments but adds no runs. Robust SSL moves quickly, and future methods may improve transfer through multi-task losses, stronger augmentations, or certified objectives. The conclusion warns against assuming transfer; it does not prove transfer impossible.
